\documentclass[border=10pt]{standalone}
\usepackage[no-math]{fontspec}
\usepackage[brazil]{babel}
\usepackage{amsmath,amssymb}
\IfFileExists{C:/Windows/Fonts/Lato_Regular.ttf}{%
  \setmainfont{Lato}[
    Path = C:/Windows/Fonts/,
    UprightFont = Lato_Regular.ttf,
    BoldFont = Lato_Bold.ttf,
    ItalicFont = Lato_Italic.ttf,
    BoldItalicFont = Lato_BoldItalic.ttf
  ]
  \setsansfont{Lato}[
    Path = C:/Windows/Fonts/,
    UprightFont = Lato_Regular.ttf,
    BoldFont = Lato_Bold.ttf,
    ItalicFont = Lato_Italic.ttf,
    BoldItalicFont = Lato_BoldItalic.ttf
  ]
}{%
  \RequirePackage[default]{lato}
}
\renewcommand{\familydefault}{\sfdefault}

\usepackage{xcolor}
\usepackage{tikz}
\usepackage{pgfplots}
\pgfplotsset{compat=1.18}
\usetikzlibrary{arrows.meta,positioning,shapes.geometric,fit,backgrounds,calc,decorations.pathreplacing,patterns}

% Paleta institucional UEFS / Extensão Rural
\definecolor{erNavy}{HTML}{1F4268}
\definecolor{erNavyDark}{HTML}{152E4A}
\definecolor{erGold}{HTML}{B38A3D}
\definecolor{erGoldDark}{HTML}{8C6927}
\definecolor{erTeal}{HTML}{0E7C7B}
\definecolor{erForest}{HTML}{2D6A4F}
\definecolor{erAlert}{HTML}{9E2A2B}
\definecolor{erInk}{HTML}{26313D}
\definecolor{erMuted}{HTML}{5B6573}

\definecolor{bgNavy}{HTML}{EDF2F7}
\definecolor{bgGold}{HTML}{F8F3E9}
\definecolor{bgTeal}{HTML}{EBF6F5}
\definecolor{bgForest}{HTML}{EEF5F0}
\definecolor{bgAlert}{HTML}{FDF2F2}

\tikzset{
  >=Stealth,
  erBox/.style={
    draw=#1, fill=white, rounded corners=5pt, line width=1.15pt,
    inner sep=7pt, align=center, font=\footnotesize\sffamily, text=erInk
  },
  erTint/.style={
    draw=#1, fill=#1!8!white, rounded corners=5pt, line width=1.15pt,
    inner sep=7pt, align=center, font=\footnotesize\sffamily, text=erInk
  },
  erBadge/.style={
    fill=#1, text=white, rounded corners=9pt, inner xsep=8pt, inner ysep=3pt,
    font=\scriptsize\bfseries\sffamily
  },
  erTitle/.style={font=\normalsize\bfseries\sffamily, text=erNavy, align=center},
  erSub/.style={font=\footnotesize\sffamily, text=erMuted, align=center},
  erArrow/.style={->, >=Stealth, line width=1.35pt},
  erDArrow/.style={<->, >=Stealth, line width=1.35pt}
}


\begin{document}
\begin{tikzpicture}[font=\sffamily]

% Título
\node[erTitle, font=\large\bfseries\sffamily] (title) at (8.80, 4.35)
  {Ciclo de Elaboração de Projetos de Extensão Rural e Marco Lógico};
\node[erSub] (sub) at (8.80, 3.88)
  {Encadeamento metodológico do diagnóstico participativo à avaliação de impacto no território};

% Pipeline horizontal de 5 etapas balanceadas (mesma altura e espaçamento amplo para setas)
\node[erTint=erTeal, text width=2.55cm, minimum height=2.15cm] (s1) at (1.60, 2.25) {
  \textbf{\textcolor{erTeal}{1. Diagnóstico}}\\
  \textbf{\textcolor{erNavy}{Participativo}}\\[5pt]
  \scriptsize Leitura da realidade\\e demandas locais
};

\node[erTint=erNavy, text width=2.55cm, minimum height=2.15cm] (s2) at (5.20, 2.25) {
  \textbf{\textcolor{erNavy}{2. Árvore de}}\\
  \textbf{\textcolor{erNavy}{Problemas \&}}\\
  \textbf{\textcolor{erNavy}{Objetivos}}\\[5pt]
  \scriptsize Causas, efeitos\\e meios-fins
};

\node[erTint=erGold, text width=2.55cm, minimum height=2.15cm] (s3) at (8.80, 2.25) {
  \textbf{\textcolor{erGoldDark}{3. Matriz de}}\\
  \textbf{\textcolor{erNavy}{Marco Lógico}}\\[5pt]
  \scriptsize Indicadores, fontes\\e pressupostos
};

\node[erTint=erForest, text width=2.55cm, minimum height=2.15cm] (s4) at (12.40, 2.25) {
  \textbf{\textcolor{erForest}{4. Cronograma,}}\\
  \textbf{\textcolor{erNavy}{Orçamento \&}}\\
  \textbf{\textcolor{erNavy}{PRONAF}}\\[5pt]
  \scriptsize Viabilidade técnica\\e financeira
};

\node[erTint=erNavy, text width=2.55cm, minimum height=2.15cm] (s5) at (16.00, 2.25) {
  \textbf{\textcolor{erNavy}{5. Monitoramento}}\\
  \textbf{\textcolor{erNavy}{\& Avaliação}}\\[5pt]
  \scriptsize Aprendizagem e\\ajuste contínuo
};

% Setas do pipeline
\draw[erArrow, color=erNavy] (s1.east) -- (s2.west);
\draw[erArrow, color=erNavy] (s2.east) -- (s3.west);
\draw[erArrow, color=erNavy] (s3.east) -- (s4.west);
\draw[erArrow, color=erNavy] (s4.east) -- (s5.west);

% Conector da Matriz de Marco Lógico para o cartão de destaque
\draw[erArrow, color=erGoldDark] (s3.south) -- (8.80, 0.50);

% Cartão inferior de destaque (erGold)
\node[draw=erGold, fill=bgGold, rounded corners=6pt, line width=1.4pt,
      text=erInk, text width=16.8cm, inner sep=10pt, align=center, anchor=north] (card) at (8.80, 0.48) {
  {\small\textbf{\textcolor{erGoldDark}{Coerência Vertical do Marco Lógico}}}\\[4pt]
  {\footnotesize\textit{Atividades $\rightarrow$ Produtos $\rightarrow$ Propósito $\rightarrow$ Financiamento Sustentável e Impacto Social}}
};

\end{tikzpicture}
\end{document}
